\documentclass[11pt]{article}
\usepackage[margin=1.0in]{geometry}
\usepackage{booktabs}
\usepackage{amsmath}
\usepackage[T1]{fontenc}
\usepackage{lmodern}
\usepackage{microtype}
\usepackage{longtable}
\usepackage{graphicx}
\usepackage[hidelinks]{hyperref}
\usepackage{underscore}

\title{Ancestry instrumentation for \texttt{fmix}\\
\large What we built to measure mixing, and everything it has told us}
\date{30 July 2026}
\author{}

\begin{document}
\maketitle

\begin{abstract}
\noindent This is a single reference for the ancestry instrumentation added to
\texttt{fmix} between 28 and 30 July 2026, and for every experiment run with it.
Part~I describes the measurements: why the previous positional meters failed,
the exact ancestor-set instrument that replaced them, the \emph{sampled} version
that removed its size ceiling, the positional-coverage measures, and the joint
generation~$\times$~ancestry census. Part~II reports the experiments and their
results, in the order they were run. Numbers here are taken from the run logs
rather than recalled; where an earlier reading has since been corrected, the
correction is stated explicitly.
\end{abstract}

\tableofcontents

\part*{Part I --- The instrumentation}
\addcontentsline{toc}{section}{Part I --- The instrumentation}

\section{The meters that failed, and why}

\texttt{fmix} carried three positional-mixing meters from the directional-walk
work: \texttt{odiff} (per-origin-family spread), \texttt{oadj} (adjacent-origin
autocorrelation) and \texttt{disp} (mean normalised displacement). All three read
a per-gate \texttt{origin} tag: the index of the input gate a gate descends from.

On the restructured menu they all reported that nothing had moved:
$\texttt{odiff}=0.0001$, $\texttt{oadj}=1.0000$ after 40M moves and $2.58\times$
growth. That reading was an artifact. A DB splice spanning \emph{mixed lineage}
stamps its products \texttt{ORIGIN\_SYNTH}, and all three meters \emph{skip}
those gates. They were therefore computed over exactly the material mixing had
failed to touch, and became more selective the better the mixing worked ---
survivorship bias pointing the wrong way.

\texttt{osyn} was added to make the bias visible: the fraction of gates whose
origin label has been destroyed. It reaches $50.7\%$ by 1M moves and $1.000$ on
the small-circuit runs, where the three meters are not merely degraded but
undefined.

\medskip
\noindent\textbf{Rule.} Read \texttt{osyn} first. Once it is high,
\texttt{odiff}/\texttt{oadj}/\texttt{disp} carry no information.

\section{Exact ancestry: sets, not labels}

A set union has no such failure mode --- merging two lineages \emph{adds}
information where a scalar label must discard it.

Each \emph{litter} (the set of gates one replacement emitted) carries the set of
original input gates that contributed to it, as a bitset. The rules are:

\begin{itemize}
\item input gate $i$ is litter $i$, with ancestor set $\{i\}$;
\item a \textbf{splice} unions its outgoing window's litters into a fresh litter;
\item a \textbf{merge} unions both parents;
\item \textbf{born-random} material (twist brackets, insert pairs) has none.
\end{itemize}

Two choices keep it affordable: singleton sets are not stored (any id below the
input count with no map entry denotes $\{id\}$, making initialisation $O(1)$
instead of $O(\text{input}^2)$), and sets are pruned to live litters at each
report, bounding the map by live litters rather than total splices.

\subsection{The three readouts}

\begin{center}
\begin{tabular}{llp{7.4cm}}
\toprule
meter & per & question \\
\midrule
\texttt{anc}    & current gate & how many original gates does this gate descend from? \\
\texttt{span}   & current gate & how far apart in the input were they (index range)? \\
\texttt{fanout} & \emph{input} gate & how many current gates carry information about it? \\
\bottomrule
\end{tabular}
\end{center}

\texttt{fanout} is redundant given \texttt{anc} and size: both sum the same
incidence relation, so $\overline{\texttt{fanout}} = \overline{\texttt{anc}}
\times \text{gates}/\text{inputs}$ exactly (verified: predicted 1018 vs reported
1015). \textbf{\texttt{span} is the one carrying independent information},
because it describes the \emph{geometry} of the ancestor sets rather than their
size.

Cost is $|\text{input}|$ bits per live litter, so the exact instrument refuses
above \textbf{20{,}000 input gates}. That ceiling is what motivated Section~3.

\section{Sampled ancestry: removing the size ceiling}

Instead of the full ancestor set per litter, trace $K$ randomly chosen input
gates --- \emph{tracers} --- and store a fixed $K$-bit mask per litter. This is
the \emph{column subsample} of the same incidence relation, so the
union-on-splice semantics carry over unchanged, while cost drops from
$|\text{input}|$ bits per litter to a constant. Enabled with
\texttt{--anc-samples K}.

Three implementation choices matter:

\begin{itemize}
\item \textbf{Tracer singletons are stored explicitly} ($K$ map entries at
init), which lets the union routine drop the implicit-singleton rule: an
untracked input gate then contributes nothing and costs nothing.
\item \textbf{All-zero masks are not stored.} Fresh litter ids always exceed the
input count, so a missing entry can never alias a singleton. Most litters carry
no tracer, so the map stays a small fraction of the circuit.
\item \textbf{Tracers are drawn from a dedicated RNG}, so the instrument cannot
perturb the walk. An exact-mode run and a sampled-mode run at the same
\texttt{--seed} follow the \emph{identical} chain, which is what makes the
calibration below a measurement rather than a coincidence.
\end{itemize}

Resume is safe without any state-format change: the tracer set is a pure
function of (input size, $K$, sample seed), all of which survive a resume, with
an assertion that catches a resume which changed $K$.

\subsection{Calibration against exact ground truth}

20k circuit, 300k moves, same seed. The identical \texttt{size} across all three
confirms the trajectory is untouched.

\begin{center}
\begin{tabular}{lrrr}
\toprule
 & size & mean \texttt{fanout}/input & mean \texttt{anc} \\
\midrule
exact                & 218{,}785 & 1850 & 169.2 \\
sampled $K=256$      & 218{,}785 & 1794 \ ($-3.0\%$) & 164.0 \ ($-3.1\%$) \\
sampled $K=64$       & 218{,}785 & 2004 \ ($+8.3\%$) & 183.2 \ ($+8.3\%$) \\
\bottomrule
\end{tabular}
\end{center}

Accuracy scales as $1/\sqrt{K}$, as expected. The estimator is
Horvitz--Thompson: each input has inclusion probability $K/m$, so the sampled sum
scaled by $m/K$ is unbiased.

\subsection{What sampled mode deliberately does not report}

In sampled mode \texttt{anc=} and \texttt{ancspan=} stay \textbf{zero}. A sampled
popcount is not \texttt{anc}, and a sampled index range is a biased \texttt{span}
(a column sample can only underestimate it). Filling those fields with sampled
values would repeat the field-meaning conflation that the \texttt{g57=}
vs \texttt{shaped=} split had to undo. Sampled runs report on a separate
\texttt{tracers:} line.

\section{Positional coverage: \texttt{cov} and \texttt{ent}}

Sampling gives something exact mode never reported: for each traced input gate,
\emph{where its descendants sit in the current circuit}.

\begin{itemize}
\item \textbf{\texttt{cov}} --- the current circuit is divided into 64 equal
contiguous slices by gate position; \texttt{cov} is the fraction of those slices
containing at least one descendant of the traced gate, averaged over tracers.
Range $1/64=0.016$ (all in one slice) to $1.0$ (present everywhere).
\item \textbf{\texttt{ent}} --- normalised entropy of the descendant counts over
those slices. \texttt{cov} says how far the influence reaches; \texttt{ent} says
how evenly. \texttt{cov} can be high while nearly all mass sits in one slice.
\item \textbf{\texttt{reach}} --- $(\max_{pos}-\min_{pos}+1)/\text{size}$. Two
tight clusters at opposite ends give $\texttt{reach}\approx 1$ but
$\texttt{cov}\approx 2/64$.
\end{itemize}

These are the security-facing quantities: they ask directly whether an adversary
can localise which part of the mixed circuit a given original gate went to.
Unlike \texttt{span}, they are natively samplable.

\subsection{A caveat that took measurement to pin down}

\texttt{cov} and \texttt{ent} are \textbf{scale-relative}: the 64 buckets are
always 64 equal fractions of the \emph{current} circuit, so the ruler rescales
with what it measures. At matched effective work ($p_{mix}=0.10$, eff $=15$):

\begin{center}
\begin{tabular}{lrrrrr}
\toprule
run & size & bucket width & blob extent & \texttt{cov} & $\texttt{cov}\times$size \\
\midrule
full 695k    & 1{,}130{,}431 & 17{,}663 & 55{,}391 & 0.049 & 55{,}391 \\
slice 1 100k &   166{,}368   &  2{,}600 & 36{,}762 & 0.221 & 36{,}762 \\
slice 2 100k &   166{,}289   &  2{,}598 & 41{,}529 & 0.250 & 41{,}529 \\
\bottomrule
\end{tabular}
\end{center}

The circuits differ $6.8\times$ in size and \texttt{cov} differs $4.5\times$, but
$\texttt{cov}\times\text{size}$ --- coverage as an absolute number of gates ---
agrees within $\sim\!1.4\times$. The predicted \texttt{cov} from blob extent
divided by bucket width reproduces the observed value to three decimals in all
three runs.

\medskip
\noindent\textbf{Consequence.} One input gate's descendants occupy a blob of
roughly \emph{fixed physical extent}, set by the mixing dynamics rather than by
circuit length. So: within one circuit size \texttt{cov} is a genuine transport
measure; across sizes compare $\texttt{cov}\times\text{size}$ instead. The same
applies to \texttt{ent}, with the extra wrinkle that a coarse ruler \emph{hides}
sub-bucket structure, systematically depressing \texttt{ent} on larger circuits.

\section{The joint generation $\times$ ancestry census}

\texttt{anc} and \texttt{dgen} had only ever been reported as separate
marginals, which cannot answer the question that matters: \emph{does re-encoding
depth buy ancestry?} The \texttt{gen-anc} line reports, per generation band, the
gate count and mean \texttt{anc} (and mean \texttt{span} in exact mode), plus the
Pearson $r$ of $(\texttt{dgen}, \texttt{anc})$.

\texttt{GEN\_FRESH} is a \texttt{u32::MAX} sentinel for born-random material, not
a large depth, so it is excluded from $r$ and reported as its own band. In
sampled mode the per-gate count is scaled by $m/K$ so bands stay comparable
across modes.

\section{Two derived quantities used throughout}

\paragraph{Total incidence.} $I = \sum_{\text{gates}} \texttt{anc}_i =
\overline{\texttt{anc}} \times \text{size}$, the gate\,$\times$\,input incidence
relation. This is the \textbf{schedule-invariant} transport measure. Per-gate
\texttt{anc} is $I/\text{size}$, which compression inflates without any
additional mixing --- a trap that cost us one wrong reading (Section~9).

\paragraph{Effective work.} Cumulative moves-per-gate,
\[
  \mathrm{eff} = \int \frac{dm}{\text{size}(m)},
\]
computed trapezoidally from the input size at move 0. A move on a small circuit
does more per-gate work than one on a large circuit, so raw move counts are not
comparable across schedules or circuit sizes; \texttt{eff} is.

\part*{Part II --- The experiments}
\addcontentsline{toc}{section}{Part II --- The experiments}

Unless stated otherwise, runs use \texttt{p\_db}$=1.0$ (so the thermostat never
fires), MIX with $s_{db}=9$, $p_{convex}=1$, $p_{mingen}=0.8$ and COMP with
$s_{db}=12$, $p_{convex}=0.5$, $p_{mingen}=0$, seed 101.

\section{Window geometry: convex sampling is the transport}

Separating the two window samplers (1M-move inhale, 2M-move exhale) gave the
single largest effect found in this campaign.

\begin{center}
\begin{tabular}{lrrr}
\toprule
sampler & final size & \texttt{anc} & \texttt{span} \\
\midrule
convex only ($p_{convex}=1$)     & 584{,}320 & \textbf{154.8} & 640 \\
contiguous only ($p_{convex}=0$) & 456{,}153 & \textbf{12.0}  & 20  \\
\bottomrule
\end{tabular}
\end{center}

Contiguous sampling is not slow, it is \emph{flat}: \texttt{anc} moves from 10.0
to 12.0 over 3M moves. The mechanism is direct --- the float counters read
$9{,}733{,}299$ floats for convex and \textbf{exactly $0$} for contiguous. With
$p_{convex}=0$ and $p_{db}=1$ no gate in the circuit ever changes position for
the entire run; the circuit is rewritten strictly in place.

This also explains the mechanism: \texttt{collect\_convex} \emph{physically
floats} gates together before replacing them, so material comes \emph{in} from
far apart and goes \emph{out} in one place. Transport happens on the way in.

\subsection{\texttt{--db-advance} rescues contiguous sampling}

Splice products are inserted as one contiguous block (all three DB modes;
placement is identical, only candidate selection differs). \texttt{--db-advance}
gives them the ballistic birth-advance that split pieces get.

\begin{center}
\begin{tabular}{lrr}
\toprule
final \texttt{anc} & advance OFF & advance ON \\
\midrule
convex     & 155 & 230 \\
contiguous & \textbf{12} & \textbf{216} \\
\bottomrule
\end{tabular}
\end{center}

Contiguous goes from $\texttt{anc}=12$ to $216$ ($18\times$) and
$\texttt{span}=20$ to $580$ ($29\times$) --- nearly to convex parity.

\section{Mode schedules A/B/C (20k circuit, 2M moves)}

Three ways to spend the same MIX/COMP budget: \textbf{A} phased (200k MIX then
1.8M COMP), \textbf{B} steady $p_{mix}=0.20$, \textbf{C} steady $p_{mix}=0.10$.

\begin{center}
\begin{tabular}{lrrrrr}
\toprule
 & final size & \texttt{anc} & \texttt{span} & \texttt{fanout} & incidence \\
\midrule
A phased      &  64{,}333 & 2649 & 3253 &  8521 & 170M \\
B 20\% MIX    & 121{,}835 & 3165 & 3820 & 19282 & 386M \\
C 10\% MIX    &  48{,}493 & 7941 & 8626 & 19255 & 385M \\
\bottomrule
\end{tabular}
\end{center}

\textbf{Steady interleave beats a single phased breath.} By total incidence ---
the schedule-invariant measure --- B $\approx$ C ($\approx$386M) $\gg$ A (170M).
A short inhale is spent before compression locks it in. B and C reach the
\emph{same} total but distribute it oppositely: C compresses into 48k gates so
each reaches 43\% of the input, B holds 122k gates at half that.

This retired an earlier claim that \texttt{span} saturates near 408; that was an
artifact of MIX-dominated large circuits. Under compression \texttt{span} climbs
past 8000.

\subsection{The effective-work rescale}

\begin{center}
\includegraphics[width=\textwidth]{../reports/ancestry_20260728/abc_vs_effwork_20260729.png}
\end{center}

Plotted against moves-per-gate rather than raw moves, the runs reach very
different effective work in the same 2M moves ($\approx$25, 33, 56) and the
transport curves \textbf{nearly collapse}: at 20 moves/gate, \texttt{anc} is
2.0k / 2.3k / 1.9k and \texttt{span} 2.6k / 3.1k / 2.7k. Most of C's apparent
lead was that a smaller circuit buys more per-gate work per move.

Residual signal survives the rescale: A leads at \emph{low} work (front-loaded
MIX), B's \texttt{span} genuinely saturates once its circuit grows large, and
\texttt{dmin} does \emph{not} collapse --- B holds $\approx$0.11 while
COMP-heavy A and C decay to $\approx$0.065. Growth keeps the circuit farther
from the \texttt{fcompress}-minimal form per unit work.

\section{Twists are expensive for ancestry}

\begin{center}
\includegraphics[width=\textwidth]{../reports/ancestry_20260728/abc_twistrate_20260729.png}
\end{center}

A/B/C re-run at swap-family twist rates 0, 0.002, 0.01 (final values):

\begin{center}
\begin{tabular}{lrrrrr}
\toprule
run & size & \texttt{anc} & \texttt{span} & \texttt{dmin} & \texttt{polf} \\
\midrule
A (no twist) &  64{,}333 & 2649 & 3253 & 0.063 & 0 \\
A tw 0.002   & 135{,}058 &  187 &  691 & 0.069 & 0.468 \\
A tw 0.01    & 247{,}382 &   69 &  393 & 0.057 & 0.496 \\
\midrule
B (no twist) & 121{,}835 & 3165 & 3820 & 0.097 & 0 \\
B tw 0.002   & 243{,}674 &  350 & 1012 & 0.087 & 0.405 \\
B tw 0.01    & 383{,}924 &   80 &  451 & 0.069 & 0.489 \\
\midrule
C (no twist) &  48{,}493 & 7941 & 8626 & 0.069 & 0 \\
C tw 0.002   & 133{,}665 &  425 & 1111 & 0.057 & 0.422 \\
C tw 0.01    & 262{,}115 &   62 &  369 & 0.040 & 0.487 \\
\bottomrule
\end{tabular}
\end{center}

Each rate step costs roughly an order of magnitude of ancestry, persists on the
effective-work axis, and \texttt{dmin} drifts \emph{down} rather than up --- the
brackets are not adding useful spelling diversity. The odometer \texttt{polf} is
already near its $\tfrac12$ ceiling at $0.002$, so \textbf{0.002 buys almost all
the polarity scrambling of 0.01 at a fraction of the ancestry cost}.

Mechanism: twist brackets are born-random (zero input ancestry), and a DB window
that gathers them unions less ancestry, so the DB keeps minting a large
low-ancestry population. Foreign (non-\texttt{comp}) material does \emph{not}
breed --- surviving foreign gates number \emph{fewer} than the brackets inserted
(ratio 0.5--0.9); the merge and COMP machinery removes some.

\subsection{Which does the damage: the CNOT brackets or the polarity flips?}

\begin{center}
\includegraphics[width=\textwidth]{../reports/ancestry_20260728/bpos_cnot_vs_polarity_20260729.png}
\end{center}

\texttt{--twist-neg-p 0} gives \emph{pure positive swaps}: the 3-CNOT brackets
are still inserted but no wire is negated. That isolates the two mechanisms.
$\texttt{polf}=0.000$ confirms the control flipped no polarities.

\medskip
\noindent\textbf{Correction.} An earlier reading of this experiment compared
endpoints and concluded the CNOT brackets dominate ($10.4\times$ against
$3.8\times$). That comparison was invalid: the no-twist arm had run to eff
$=32.9$ while the pure-swap arm reached only $17.4$. At \emph{matched} effective
work the two mechanisms are comparable:

\begin{center}
\begin{tabular}{lrrrrrr}
\toprule
eff & no twist & pure swap & full family & CNOTs cost & flips cost & total \\
\midrule
10.0 & 703  & 203 & 39 & $3.47\times$ & $5.20\times$ & $18.1\times$ \\
14.0 & 1332 & 276 & 66 & $4.83\times$ & $4.17\times$ & $20.2\times$ \\
16.0 & 1672 & 293 & 79 & $5.70\times$ & $3.72\times$ & $21.2\times$ \\
\bottomrule
\end{tabular}
\end{center}

So the foreign CNOT brackets \emph{do} block mixing substantially on their own
($3.5$--$5.7\times$), which was the hypothesis under test --- but the polarity
flips cost a comparable amount ($3.7$--$5.2\times$), and the split is roughly
even rather than lopsided. Turning off the negations recovers only part of the
cost; the brackets are the structural price of any swap-family twist.

\section{Generation versus ancestry}

\begin{center}
\includegraphics[width=\textwidth]{../reports/ancestry_20260728/genanc_20260730.png}
\end{center}

Retrace of A/B/C with exact ancestry and the joint census.

\begin{center}
\begin{tabular}{lrrrr}
\toprule
band & \multicolumn{1}{c}{A} & \multicolumn{1}{c}{B} & \multicolumn{1}{c}{C} \\
\midrule
g0      & 1.0  & 1.0  & 1.0 \\
g17--32 & 887  & ---  & --- \\
g33--64 & 2580 & 2027 & --- \\
g65+    & 2613 & 3310 & \textbf{7460} \\
\bottomrule
\end{tabular}
\end{center}

\textbf{Depth does buy ancestry}, monotonically. But \textbf{the same nominal
depth means very different mixing}: at band g65+, A reaches 2613 and C reaches
7460, a $2.9\times$ spread. A \emph{saturates} --- its last two bands are
$2580 \to 2613$ ($+1.3\%$) --- while B still climbs $+63\%$ across the same step.
Generation is not a protocol-independent measure of mixing.

The Pearson $r$ is the weaker instrument and should not be leaned on: it starts
at 0.47--0.71 and decays to $\approx 0$, with B going to $-0.41$. The band
populations show why --- by the end, B and C have collapsed almost entirely into
a single band, leaving no depth variance to correlate against. This is range
restriction, not a loss of the underlying relationship.

\paragraph{Span versus global generation.} The protocols do not share one curve:
at $G\approx 50$ the spans are A $\approx$ 3155, B $\approx$ 3500, C $\approx$
4300, and C ends at $G=105$/span 8107 against A at $G=51$/span 3155. C buys
$\approx 25\%$ more span per unit of global generation \emph{and} reaches twice
the generation count.

\section{Production scale: the 695k gadget}

First use of sampled ancestry where exact ancestry is impossible. Input: a fresh
Gray-code gadgetized 128-bit sliced sandwich, 695{,}075 gates / 512 wires, seed
730, self-verify passed. Three arms at $p_{mix}=0.20/0.10/0.05$, 20M moves,
$K=256$ tracers.

\begin{center}
\begin{tabular}{lrrrrrr}
\toprule
arm & $p_{mix}$ & final size & growth & eff & \texttt{anc} & \texttt{dmin} \\
\midrule
big20 & 0.20 & 1{,}884{,}758 & $2.71\times$ & 15.0 & 17{,}532 & 0.081 \\
big10 & 0.10 & 1{,}262{,}860 & $1.82\times$ & 19.8 & 26{,}678 & 0.048 \\
big05 & 0.05 &   948{,}984   & $1.37\times$ & 23.7 & 35{,}173 & 0.030 \\
\bottomrule
\end{tabular}
\end{center}

All three were still growing monotonically at 20M moves (peak size = final size
in every case). Less MIX buys a smaller circuit \emph{and} more ancestry per
gate --- the same total transport concentrated into fewer gates --- at the cost
of \texttt{dmin}, which drains from 0.081 to 0.030 as the schedule becomes more
compression-heavy.

\section{Can slices replace the full circuit?}

\begin{center}
\includegraphics[width=\textwidth]{../reports/ancestry_20260728/three_slices_20260730.png}
\end{center}

Three disjoint 100k-gate slices of the same gadget (gates 1--100k, 100--200k,
200--300k), each at the same three $p_{mix}$ values, with moves chosen so
moves-per-gate and report count match the 20M full run exactly.

\begin{center}
\begin{tabular}{lrrrrr}
\toprule
metric ($p_{mix}=0.10$) & slice 1 & slice 2 & slice 3 & s3 vs s2 & s2 vs s1 \\
\midrule
\texttt{anc}   & 22{,}774 & 29{,}009 & 30{,}035 & $+3.5\%$ & $+27.4\%$ \\
\texttt{desc}  & 40{,}673 & 51{,}965 & 53{,}695 & $+3.3\%$ & $+27.8\%$ \\
\texttt{cov}   & 0.298 & 0.363 & 0.378 & $+4.1\%$ & $+21.8\%$ \\
\texttt{dmin}  & 0.053 & 0.063 & 0.063 & $0.0\%$  & $+18.9\%$ \\
size           & 178{,}592 & 179{,}136 & 178{,}772 & $-0.2\%$ & $+0.3\%$ \\
\bottomrule
\end{tabular}
\end{center}

\textbf{Slice 1 is the outlier.} Slices 2 and 3 agree within 3--6\% on every
ancestry metric while slice 1 sits 10--54\% below both, consistently and in the
same direction across all three schedules. The gadget's opening 100k gates
--- which contain the zero-slice preblock --- genuinely mix worse. Growth is
identical across all four ($<4\%$), so tracking size alone would have wrongly
concluded the slices were equivalent; it took the ancestry instrument to see it.

Against the full circuit:

\begin{center}
\begin{tabular}{lrr}
\toprule
$p_{mix}$ & BIG vs slice 1 & BIG vs mid-slices (2,3) \\
\midrule
0.20 & $+35.3\%$ & $-8.6\%$ \\
0.10 & $+15.2\%$ & $-10.3\%$ \\
0.05 & $+14.9\%$ & $+2.1\%$ \\
\bottomrule
\end{tabular}
\end{center}

The full circuit tracks the \emph{middle} slices to within 2--10\%, as it should
--- the atypical opening is only $1/7$ of it. \textbf{Practical rule: a
mid-circuit 100k slice reproduces full-gadget behaviour to within $\sim$10\% at
$1/7$ the moves, provided the slice does not start at gate 0, and provided two
mid slices are used to bound the residual variability.}

\subsection{Validity window}

\begin{center}
\includegraphics[width=\textwidth]{../reports/ancestry_20260728/anc_progression_20260730.png}
\end{center}

Absolute \texttt{anc} is comparable between slice and full circuit only while
both are far from saturation. At the end of these runs the slices reach 18--35\%
of their 100k input while the full circuit is at $\approx$5\% of its 695k input.
Push much further and a slice will run out of distinct input gates while the
full circuit keeps going, and the agreement will break down.

\section{Reconstruction heatmaps of the mixed outputs}

Affine (degree 1) and degree-2 predictor maps for the three 20M-move outputs,
against two references: (a) the original circuit $C$, and (b) the gadgetized
circuit that \texttt{fmix} was given as input. Read by the ridge measure
(depth / $\rho$ / perm-$z$ / contrast), never by the mean.

\begin{center}
\begin{tabular}{lrrrrr}
\toprule
map & mean $H$ & depth & $\rho$ & perm $z$ & contrast \\
\midrule
(a) vs $C$, degree 1      & 0.494--0.498 & 0.047 & 0.42--0.44 & 7.7 & 0.34--0.46 \\
(a) vs $C$, degree 2      & 0.490--0.495 & 0.069 & 0.43       & 1.9 & 0.43--0.45 \\
(b) vs gadget, degree 1   & 0.465--0.470 & 0.311--0.321 & \textbf{1.000} & 10.0 & 2.92--3.22 \\
(b) vs gadget, degree 2   & 0.465--0.468 & 0.253--0.281 & \textbf{1.000} & 4.6  & 1.54--1.78 \\
\bottomrule
\end{tabular}
\end{center}

\begin{center}
\includegraphics[width=0.49\textwidth]{../reports/ancestry_20260728/hmA_d1.png}
\includegraphics[width=0.49\textwidth]{../reports/ancestry_20260728/hmB_d1.png}
\end{center}

\paragraph{(a) The logical computation is hidden.} Reading row structure rather
than $\rho$: the plate is uniformly 0.5 across every interior row. The only live
cells are the two trivial corners --- the top-left block, where $C$'s first
prefixes are recoverable because $x$ is still exposed on wires 0--127 at the
start, and a one-cell sliver at the final output. The apparent ridge at row
$\approx 0.65$ is the argmin wandering in a dead plate, which is exactly what
produces $\rho=0.42$; perm-$z$ falls to 1.9 ($p=0.057$) at degree 2. Degree 2
finds essentially nothing affine did not.

\paragraph{(b) The correspondence to its own input survives intact.} A bright
unbroken ridge corner to corner, $\rho$ exactly 1.000 in all six maps, alive in
every interior row. Given the gadget, any prefix of the mixed circuit can be
aligned to the gadget's matching prefix. This is the security-relevant readout of
the locality the tracers already measured: the mixing rewrites in place and does
not relocate material globally, so progress through the circuit stays legible.

The trend runs the way the ancestry data predicts: depth and contrast increase
monotonically as $p_{mix}$ falls ($2.92 \to 3.03 \to 3.22\sigma$).
\texttt{big05}, the arm with the \emph{most} ancestry per gate, has the
\emph{strongest} diagonal. \textbf{Ancestry mixing and positional mixing are
different axes}; compression packs more lineage into each gate without moving
anything.

\medskip
\noindent Degree-2 caveats: the grid is much coarser ($21\times21$ against
$101\times151$), so depth and contrast are not directly comparable across
degrees; and products were restricted to wires 0--63 (2529 regressors against
9216 training samples), which makes degree 2 a \textbf{lower bound} --- products
among the excluded 448 wires are unseen.

\section{What this adds up to}

\begin{enumerate}
\item \textbf{Report incidence, not \texttt{anc} alone}, when comparing
schedules or sizes. \texttt{anc} per gate is $I/\text{size}$, and compression
inflates it without additional mixing.
\item \textbf{Convex sampling is the transport mechanism}; contiguous sampling
moves nothing at all unless \texttt{--db-advance} is on. The historical
$p_{convex}=0.5$ default was discarding half the transport.
\item \textbf{Steady MIX/COMP interleave beats a phased breath} at equal budget.
Whether \emph{cyclic} breathing under a size band beats steady interleave ---
C-like span at B-like \texttt{dmin} within a bounded footprint --- remains the
open test.
\item \textbf{Twists cost roughly an order of magnitude of ancestry per rate
step}, split roughly evenly between the foreign CNOT brackets and the polarity
flips. Use the lowest rate that moves \texttt{polf}; 0.002 already nearly
saturates it.
\item \textbf{Depth is not a protocol-independent measure of mixing.} The same
generation band carries $2.9\times$ different ancestry under different schedules.
\item \textbf{Mid-circuit slices are a valid $7\times$ cheaper proxy} for
schedule comparison, but the gadget's opening region is not representative.
\item \textbf{The mixed circuit hides $C$ but not its own input.} Against $C$ the
reconstruction plate is dead to both affine and degree-2 adversaries; against the
gadget the diagonal is perfect. Closing that gap requires global relocation of
material, which is precisely what the positional measures say is not happening.
\end{enumerate}

\section{Open}

\begin{itemize}
\item A degree-2 map with the product-wire set straddling encoding blocks rather
than taking a prefix, to strengthen the (a) statement.
\item Cyclic breathing under a size band, against steady interleave.
\item An absolute-scale positional measure (descendant spread in gates) so
\texttt{cov}/\texttt{ent} become size-comparable without the
$\texttt{cov}\times\text{size}$ correction.
\item In flight at the time of writing: a nine-arm battery on slice 2 --- three
g57-mode twist rates at $p_{mix}=0.05$, and a six-point $p_{mix}$ sweep from
0.01 to 0.07 --- all to 15 effective moves per gate.
\end{itemize}

\end{document}
